%%%%%%%%%%%%%%%%%%%%%%%%%%%%%%%%%%%%%%%%%%%%%%%%%%%%%%%%%%%%%%%%%%%%%%
% Lớp 10 | Phương pháp tọa độ trong mặt phẳng (𝑂𝑥𝑦) | Đường tròn trong mặt phẳng tọa độ
% Dạng 0H9?4-4: Vị trí tương đối liên quan đường tròn
% Số câu: 11  (sinh tự động bởi tach_cau.py; từ nay thầy sửa trực tiếp file này)
%%%%%%%%%%%%%%%%%%%%%%%%%%%%%%%%%%%%%%%%%%%%%%%%%%%%%%%%%%%%%%%%%%%%%%

%%%%%%%%%%%%%%%% TRẮC NGHIỆM 4 PHƯƠNG ÁN %%%%%%%%%%%%%%%%

% nguồn: MaHoaNH_TracNghiem_Lop10
\begin{ex}%[0H9H4-4]%[ma:d39ea1c6]
	Trong mặt phẳng $Oxy$, đường tròn $(C)$ có tâm $I(-1;3)$ và tiếp  xúc với đường thẳng $d\colon 3x-4y+5=0$ tại điểm $H$ có tọa độ là
	\choice 
		{$\left(-\dfrac{1}{5};-\dfrac{7}{5} \right)$} 
		 {\True $\left(\dfrac{1}{5};\dfrac{7}{5} \right)$}
		 {$\left(\dfrac{1}{5};-\dfrac{7}{5} \right)$}
		 {$\left(-\dfrac{1}{5};\dfrac{7}{5} \right)$}
	\loigiai{ Ta  có
		$IH  \perp d\Rightarrow IH: 4x+3y+c=0$. \\
		 Đường thẳng $IH$ qua $I(-1; 3)$ nên $4(-1)+3\cdot 3+c=0\Leftrightarrow c=-5$.\\
		Vậy  $IH\colon 4x+3y-5=0$.\\
		Giải hệ $\heva{&4x+3y-5=0\\&3x-4y+5=0}  \Leftrightarrow \heva{&x=\dfrac{1}{5}\\&y=\dfrac{7}{5}} \Rightarrow H\left(\dfrac{1}{5}; \dfrac{7}{5} \right)$.
	 }
\end{ex}

% nguồn: MaHoaNH_TracNghiem_Lop10
\begin{ex}%[0H9H4-4]%[ma:b1c322ca]
	Trong mặt phẳng $Oxy$, tọa độ giao điểm của đường thẳng $\Delta\colon x-2y+3=0$ và đường tròn $(C)\colon x^2+y^2-2x-4y=0$ là 
	 \choice
		{\True $(3;3)$ và  $(-1;1)$}
		{$(-1;1)$ và $(3;-3)$} 
		{$(3;3)$ và  $(1;1)$}
		{$(2;1)$ và  $(2;-1)$}
	\loigiai{ 
		Tọa độ giao điểm là nghiệm của hệ phương  trình sau
		\begin{eqnarray*}
			 & & \heva{ & x-2y+3=0 \\
			&  x^2+y^2-2x-4y=0
			 } \Leftrightarrow
			\heva{ & x=2y-3 \\ 
			 & (2y-3)^2+y^2-2(2y-3)-4y=0
			} \\ 
			& \Leftrightarrow & \heva{&{y^2-4y+3=0} \\ & x=2y-3} \Leftrightarrow \hoac{& \heva{& y=1 \\ & x=-1}  \\& \heva{& y=3 \\ & x=3.}}
		\end{eqnarray*}
		Vậy tọa độ giao điểm là $(3;3)$ và  $(-1;1)$.
	 }
\end{ex}

% nguồn: MaHoaNH_TracNghiem_Lop10
\begin{ex}%[0H9H4-4]%[ma:fad7de3d]
	Trong mặt phẳng $Oxy$, cho  đường tròn $(C)\colon x^2+y^2-4x+6y-3=0$. Xét các mệnh đề sau
	\begin{itemize}
		\item[(1)] Điểm $A(1;1)$ nằm ngoài  $(C)$.
		\item[(2)] Điểm $O(0;0)$ nằm trong  $(C)$.
		\item[(3)] $(C)$ cắt trục  tung tại hai điểm phân biệt.
	\end{itemize}
	Mệnh  đề đúng là
	\choice 
		{Chỉ (1)} 
		 {Chỉ (2)}
		{Chỉ (3)} 
		{\True Cả (1), (2) và  (3)}
	 \loigiai{
		Đặt $f(x; y)=x^2+y^2-4x+6y-3$.\\ 
		$f(1; 1)=1+1-4+6-3=1> 0\Rightarrow A$ ở ngoài $(C)$.\\ 
		$f(0; 0)=-3< 0\Rightarrow O(0;  0)$ ở trong $(C)$.\\
		$x=0\Rightarrow y^2+6y-3=0$.\\ 
		Phương trình này có hai nghiệm, suy  ra $(C)$ cắt $y'Oy$ tại $2$ điểm.
	 }
\end{ex}

% nguồn: MaHoaNH_TracNghiem_Lop10
\begin{ex}%[0H9V4-4]%[ma:e39f6983]
	Trong mặt phẳng $Oxy$, cho  đường tròn $(C)\colon x^2+y^2-4x-6y+5=0$. Đường thẳng $d$ đi qua $A(3;2)$ và cắt $(C)$ theo một dây cung ngắn nhất có phương trình là
	\choice 
		{$2x-y+2=0$} 
		 {$x+y-1=0$}
		 {\True $x-y-1=0$}
		{$x-y+1=0$} 
	\loigiai{ 
		\begin{center}
			\begin{tikzpicture}[line join=round,line cap=round,thick,scale= 0.9]
					 \def\r{3}
					\def\goc{65} 
					 \def\gochai{170}
					 \path
						(0,0) coordinate (I) 
						++(\goc:\r)  coordinate (N)
						 (I)++(\gochai:\r) coordinate (M)
						 ($(N)!0.5!(M)$) coordinate (H)
						($(H)!0.65!(M)$) coordinate  (A);
					\foreach \x/\y/\z  in {
						 I/H/M}{\draw pic[draw,angle radius=2.5mm]{right angle=\x--\y--\z};}
					\draw (M)--(N) 
						(I)--(H) 
						(I)--(A); 
					\draw  [blue]let \p1 = ($(I) - (M)$) in (I) circle ({veclen (\x1,\y1)});
					\foreach \x/\y  in {
						I/-90,M/180,A/85,H/85,N/30 
						}{\draw[fill=black] (\x) circle (1pt)++(\y:0.35) node  {$\x$};}
				\end{tikzpicture}
		\end{center}
		Ta có $f(x;y)=x^2+y^2-4x-6y+5\Rightarrow f(3; 2)=9+4-12-12+5=-6  < 0$.\\
		Vậy $A(3; 2)$ ở trong $(C)$.\\ 
		Dây cung $MN$ ngắn  nhất $\Leftrightarrow IH$ lớn nhất $\Leftrightarrow H\equiv A$ $\Leftrightarrow MN$ có véc-tơ pháp tuyến là $\overrightarrow{IA}=(1; -1)$. \\
		Vậy $d$ có phương trình là $1(x-3)-1(y-2)=0\Leftrightarrow  x-y-1=0$.
	 }
\end{ex}

% nguồn: MaHoaNH_TracNghiem_Lop10
\begin{ex}%[0H9V4-4]%[ma:c7587191]
	Trong mặt phẳng $Oxy$, cho đường tròn $(C)\colon x^2+y^2+6x-2y+5=0$ và đường thẳng $d$ đi qua điểm $A(-4;2)$, cắt $(C)$ tại hai điểm $M$, $N$ sao cho $A$ là trung điểm của  $MN$. Phương trình của đường thẳng $d$ là
	 \choice
		 {\True $x-y+6=0$}
		{$7x-3y+34=0$} 
		 {$7x-3y+30=0$}
		 {$7x-y+35=0$}
	 \loigiai{
	 \immini{
		$(C)$ có tâm $I(-3;1),R=\sqrt{5}$. Do đó $IA=\sqrt{2} < R \Rightarrow A$ ở trong $(C)$.\\ 
		$A$ là trung điểm của $MN\Rightarrow IA\perp MN\Rightarrow \overrightarrow{IA}=(-1;1)$ là vectơ  pháp tuyến của $d$, nên $d$ có phương trình là
		\[ 
		 -1(x+4)+1(y-2)=0\Leftrightarrow x-y+6=0.
		\] 
		}{ 
		\begin{tikzpicture}[line join=round,line cap=round,thick,scale= 0.6]
				\def\r{3} 
				\def\goc{65} 
				 \def\gochai{170}
				\path 
					(0,0) coordinate  (I)
					++(\goc:\r) coordinate (N) 
					(I)++(\gochai:\r) coordinate  (M)
					($(N)!0.5!(M)$) coordinate  (A);
				 \foreach \x/\y/\z in {
					I/A/M}{\draw pic[draw,angle  radius=2.5mm]{right angle=\x--\y--\z};}
				 \draw (M)--(N)
					(I)--(A); 
				\draw [blue]let  \p1 = ($(I) - (M)$) in (I) circle ({veclen (\x1,\y1)});
				\foreach  \i/\g in {
					I/-90,M/180,A/95,N/30 
					 }{\draw[fill=white](\i) circle (1pt)
					($(\i)+(\g:3mm)$) node[scale=0.7]{$\i$};} 
			\end{tikzpicture}
		 }
	 }
\end{ex}

% nguồn: MaHoaNH_TracNghiem_Lop10
\begin{ex}%[0H9V4-4]%[ma:3c0996cc]
	Trong mặt phẳng $Oxy$, cho đường tròn $(C)\colon x^2+y^2-2x+6y+6=0$ và đường thẳng $d \colon  4x-3y+5=0$. Đường thẳng $d'$ song song với đường thẳng $d$ và chắn trên $(C)$ một dây cung có độ dài bằng $2\sqrt{3}$ có phương trình là
	 \choice
		{$4x-3y+8=0$} 
		{\True $4x-3y-8=0$  hoặc $4x-3y-18$}
		{$4x-3y-8=0$} 
		{$4x+3y+8=0$} 
	 \loigiai{
	\immini{ 
		$(C)$  có tâm $I(1; -3)$ và bán kính $R=2$.\\
		 $d' \parallel d\Rightarrow d'$ có phương trình $4x-3y+m=0$, $\left(m\ne 5\right)$.\\
		Vẽ $IH\perp  MN\Rightarrow HM=\sqrt{3} \Rightarrow IH^2=R^2-HM^2=4-3=1$.\\
		Ta  có $ \mathrm{d}\left(I,d'\right)=IH\Leftrightarrow \dfrac{\left|4\cdot 1-3\cdot (-3)+m\right|}{\sqrt{16+9}}=1$\\
		$\Leftrightarrow \left|m+13\right|=5\Leftrightarrow \hoac{& m=-8 \\ &  m=-18.}$\\
		Vậy $d'\colon 4x-3y-8=0$ hoặc  $d'\colon 4x-3y-18=0$.
		}{ 
		\begin{tikzpicture}[line join=round,line cap=round,thick,scale= 0.6]
				\def\r{3} 
				 \def\goc{65}
				 \def\gochai{170}
				\path 
					 (0,0) coordinate (I)
					 ++(\goc:\r) coordinate (N)
					(I)++(\gochai:\r) coordinate  (M)
					($(N)!0.5!(M)$) coordinate  (H);
				\foreach \x/\y/\z in  {
					I/H/M}{\draw pic[draw,angle radius=2.5mm]{right angle=\x--\y--\z};} 
				 \draw (M)--(N)
					(I)--(H); 
				\draw [blue]let \p1 = ($(I) - (M)$) in  (I) circle ({veclen (\x1,\y1)});
				\foreach \i/\g in { 
					 I/-90,M/180,H/95,N/30
					}{\draw[fill=white](\i) circle (1pt) 
					($(\i)+(\g:3mm)$)  node[scale=0.7]{$\i$};}
			\end{tikzpicture}
		} 
	} 
\end{ex}


%%%%%%%%%%%%%%%% ĐÚNG/SAI %%%%%%%%%%%%%%%%

% nguồn: MaHoaNH_DungSai_Lop10
\begin{ex}%[0H9N4-4]%[ma:6f8a3aaa]
	Cho đường  tròn $(C)\colon x^2+(y-2)^2=4$. Các khẳng định sau đây, khẳng định nào đúng, khẳng định nào sai?
	\choiceTF 
		{\True Điểm $A(2;1)$  nằm ngoài đường tròn $(C)$}
		{Điểm $B(1;1)$  nằm trên đường tròn $(C)$}
		{\True Đường thẳng $d\colon  y-4=0$ tiếp xúc với $(C)$}
		{\True  Đường thẳng $\Delta \colon x+2y-3=0$ cắt $(C)$ tại hai điểm phân biệt}
	 \loigiai{
		Đường tròn $(C)$ có tâm $I(0;2)$, bán kính $R=2$. 
		\begin{itemchoice}
			\itemch Ta có $IA=\sqrt{(2-0)^2+(1-2)^2}=\sqrt{5}>R$ nên $A$ nằm ngoài đường tròn  $(C)$.
			\itemch  Ta có $IB=\sqrt{(1-0)^2+(1-2)^2}=\sqrt{2}<R$ nên điểm $B$ nằm bên trong đường tròn $(C)$.
			\itemch Ta có $\mathrm{d}(I,d)=\dfrac{|2-4|}{\sqrt{0^2+1^2}}=2=R$ nên $d$ tiếp xúc với $(C)$. 
			\itemch Ta có $\mathrm{d}(I,\Delta)=\dfrac{|0+4-3|}{\sqrt{1^2+2^2}}=\dfrac{1}{\sqrt{5}}<R$  nên $\Delta $ cắt $(C)$ tại hai điểm phân biệt.
		\end{itemchoice}
	 }
\end{ex}

% nguồn: MaHoaNH_DungSai_Lop10
\begin{ex}%[0H9H4-4]%[ma:3e8dfe0d]
	Trong hệ trục tọa độ $Oxy$, cho đường tròn $(C)\colon (x+1)^2+(y-2)^2=9$. Xét tính đúng - sai của các  khẳng định bên dưới.
	 \choiceTF
		{\True  $(C)$ có bán kính $R=3$}
		{$(C)$ có tâm  là $I(1;-2)$}
		{\True $(C)$ cắt $Ox$  tại hai điểm phân biệt}
		{$(C)$ tiếp xúc với đường  thẳng $d\colon 3x-4y+1=0$}
	\loigiai{Đường tròn $(C)\colon (x+1)^2+(y-2)^2=9$ có tâm  $I(-1;2)$, bán kính $R=3$.
		\begin{itemchoice}
			\itemch  $(C)$ có bán kính $R=3$ là đúng.
			\itemch  $(C)$ có tâm là $I(1;-2)$ là sai.
			\itemch Ta có  $\mathrm{d}(I,Ox)=\left|y_I\right|=2<R$ nên $(C)$ cắt $Ox$ tại hai điểm phân biệt là đúng.
			\itemch Ta  có $\mathrm{d}(I,d)=\dfrac{|-3-8+1|}{\sqrt{3^2+(-4)^2}}=2<R$ nên $(C)$ tiếp xúc với $d$ là sai.
	\end{itemchoice}}
\end{ex}

% nguồn: MaHoaNH_DungSai_Lop10
\begin{ex}%[0H9H4-4]%[ma:278f996c]
	Các khẳng định sau  đây, khẳng định nào đúng, khẳng định nào sai?
	 \choiceTF
		{\True  Điểm $B(-1;-2)$ nằm trên đường tròn $(C)\colon (x-1)^2+(y+1)^2=5$}
		{\True Nếu khoảng cách từ tâm đường tròn $(C)$ đến đường thẳng $d$ nhỏ hơn bán kính thì đường thẳng $d$ cắt đường tròn $(C)$ tại hai điểm  phân biệt}
		{\True Điểm $M(1;3)$ là giao điểm của đường thẳng $d\colon x-y+2=0$  và đường tròn $(C)\colon (x-2)^2+(y-1)^2=5$}
		{\True Đường thẳng $\Delta$ đi qua điểm $A(-2;1)$ luôn cắt đường tròn $(C)\colon (x+1)^2+(y-1)^2=4$ tại hai điểm  phân biệt}
	 \loigiai{
		\begin{itemchoice}
			\itemch Vì $(-1-1)^2+(-2+1)^2=5$ nên điểm $B(-1;-2)$ nằm trên đường  tròn $(C)\colon (x-1)^2+(y+1)^2=5$.
			\itemch Nếu khoảng cách từ tâm đường tròn $(C)$ đến đường thẳng $d$ nhỏ hơn bán kính thì đường thẳng $d$ cắt  đường tròn $(C)$ tại hai điểm phân biệt.
			\itemch Ta thấy điểm $M(1;3)$ thuộc đường thẳng $d\colon x-y+2=0$ và cũng thuộc đường tròn $(C)\colon (x-2)^2+(y-1)^2=5$ nên $M$ là giao điểm của  $d$ và $(C)$.
			\itemch Đường  tròn $(C)$ có tâm $I(-1;1)$, bán kính $R=2$ và $IA=\sqrt{(-2+1)^2+(1-1)^2}=1$.\\
			Đường thẳng $\Delta$ đi qua điểm $A$ nên $\mathrm{d}(I,\Delta)\leq IA<R$ nên  đường thẳng $\Delta$ luôn cắt đường tròn $(C)$ tại hai điểm phân biệt.
		\end{itemchoice}
	} 
\end{ex}

% nguồn: MaHoaNH_DungSai_Lop10
\begin{ex}%[0H9H4-4]%[ma:768de62d]
	Trong mặt phẳng $Oxy$, cho đường tròn $(C)$ đi qua ba điểm $A(2;0)$, $B(0; -3)$, $C(5;-3)$ tâm $I$. 
	\choiceTF 
		{\True  Đường kính của đường tròn $(C)$ bằng $\sqrt{26}$}
		{\True Tung độ của $I$  là $-\dfrac{5}{2}$}
		{\True  Đường tròn $(C)$ đi qua điểm $N(3;0)$}
		 {$IO=5\sqrt{2}$}
	 \loigiai{
		Gọi tâm  đường tròn là $I (a;b)$. Theo giả thiết $\heva{&AI^2=BI^2\\
		 &AI^2=CI^2}$\\
		\[\Leftrightarrow \heva{&(a-2)^2+b^2=a^2+(b+3)^2\\ 
		&(a-2)^2+b^2=(a-5)^2+(b+3)^2} 
		\Leftrightarrow \heva{&4a+6b=-5\\ 
		 &6a-6b=30}
		 \Leftrightarrow \heva{&a=\dfrac{5}{2} \\
		 &b=-\dfrac{5}{2}.}
		 \]
		Bán kính đường tròn là $R=\sqrt{\left(\dfrac{5}{2}-2\right)^2+\left(-\dfrac{5}{2} \right)^2}=\sqrt{\dfrac{13}{2}}$.\\ 
		Vậy phương trình đường tròn $(C)\colon \left(x-\dfrac{5}{2} \right)^2+\left(y+\dfrac{5}{2} \right)^2=\dfrac{13}{2}$.\\ 
		Khi  đó
		\begin{itemchoice}
			 \itemch Đường kính của đường tròn $(C)$ bằng $\sqrt{26}$.
			\itemch  Tung độ của $I$ là $-\dfrac{5}{2}$.
			\itemch Đường tròn $(C)$ đi qua điểm  $N(3;0)$.
			\itemch $IO=\dfrac{5\sqrt{2}}{2}$. 
		\end{itemchoice}
	} 
\end{ex}


%%%%%%%%%%%%%%%% TRẢ LỜI NGẮN %%%%%%%%%%%%%%%%

% nguồn: MaHoaNH_TraLoiNgan_Lop10
\begin{bt}%[0H9H4-4]%[ma:958fcb7c]
	Trong  mặt phẳng $Oxy$, cho đường tròn $(C)$ có tâm $B(1;1)$ và cắt $d \colon 3x+4y+8=0$ tại $M$, $N$ thoả mãn $MN=8$. Bán kính của $(C)$ bằng bao nhiêu?
	\shortans{5} 
	 \loigiai{
	\immini{ 
		Gọi $H$ là hình  chiếu của $B$ lên $d\colon 3x+4y+8=0$. Khi đó khoảng cách từ điểm $B$ đến đường thẳng $d$ là
		 \[
		 BH=\dfrac{|3\cdot 1+4\cdot 1+8|}{\sqrt{3^2+4^2}}=3.
		 \]
		$H$ là trung điểm của  $MN$ nên $HM=4$. Suy ra bán kính đường tròn $(C)$ là $R=\sqrt{BH^2+HM^2}=\sqrt{3^2+4^2}=5$.
		}{ 
		\begin{tikzpicture}[line join=round,line cap=round,thick,scale= 0.6]
				\def\r{3} 
				 \def\goc{65}
				 \def\gochai{170}
				\path 
					(0,0) coordinate (B) 
					++(\goc:\r)  coordinate (N)
					(B)++(\gochai:\r) coordinate  (M)
					($(N)!0.5!(M)$) coordinate (H); 
				\foreach \x/\y/\z  in {
					 B/H/M}{\draw pic[draw,angle radius=2.5mm]{right angle=\x--\y--\z};}
				 \draw (M)--(N)
					 (B)--(H);
				\draw [blue]let \p1 = ($(B) -  (M)$) in (B) circle ({veclen (\x1,\y1)});
				 \foreach \i/\g in {
					B/-90,M/180,H/95,N/30 
					}{\draw[fill=white](\i) circle  (1pt)
					 ($(\i)+(\g:3mm)$) node[scale=0.7]{$\i$};}
			\end{tikzpicture}
		 }
	 }
\end{bt}
